

Pafnuty Chebychev was a master integrator. In 1853 he made several advances to the theory of algebraic integration. In particular, he proved the following theorem which has the (possibly unique) distinction of appearing on a Russian postage stamp in 2021\cite{chebychev_stamp}. 

\begin{figure}[H]
\centering
\includegraphics[width=0.33\textwidth]{Pafnuty_Chebyshev_2021_stamp_of_Russia.jpg}
\end{figure}
 
\medskip

\noindent\textsf{Chebychev's Theorem} \qquad {\it If $p$, $q$, $r$ are rational numbers and $a$, $b$ are real numbers with $a,b,r \neq 0$, then
$$\int x^p \left(a\,x^r+b\right)^q,$$
is integrable in terms of elementary functions if and only if at least one of
$$q,\frac{p+1}{r}, \text{ or } q+\frac{p+1}{r}$$
is an integer.}

\medskip 

We will summarise the sufficiency case where each of $q\in\mathbb{Z}$, $(p+1)/r\in\mathbb{Z}$, and $q+(p+1)/r\in\mathbb{Z}$ produces elementary antiderivatives. For the proof of the necessity case of Chebychev's Theorem is beyond the scope of this book and we refer the reader to Ritt\cite[pp. 37]{ritt1948}.\\

\noindent\textsc{Chebychev type I} \qquad Where $q \in \mathbb{Z}$, we substitute $u = x^{1/N}$, where $N = \text{lcm}(\text{den}(p), \text{den}(r))$ and obtain a polynomial integrand in $u$. \\

\noindent\textsc{Chebychev type II} \qquad Where $(p+1)/r \in \mathbb{Z}$, we substitute $u = \left(a\,x^r+b\right)^N$, where $N=1/\text{den}(q)$ and obtain a rational integrand in $u$. \\

\noindent\textsc{Chebychev type III} \qquad Where $q+(p+1)/r \in \mathbb{Z}$, we substitute $u=x^r$, then $t = \left((a\,u+b)/u\right)^N$, where $N = 1/\text{den}(q)$ and we obtain a rational integrand in $t$.\footnote{Thank you to Dia Jivan, a high school student from the UK for correcting an error in the explanation of this case.}\\

For example, consider the following integral 
$$\int x^{-5/6}\sqrt{\sqrt[3]{x} - 1}\, dx.$$
We have $p = -5/6$, $q = 1/2$, and $r = 1/3$. Then as $q + (p + 1)/r = 1$, this is a Chebychev type III integral. Let $u = \sqrt[3]{x}$, then $du = 1/3 x^{-2/3} dx$ and we have 
$$\int \frac{3\sqrt{u-1}}{\sqrt{u}}du.$$
In order to make the substitution $t = \left((u-1)/u\right)^N$, where $N = 1/\text{den}(q)$, we must write the integral as 
$$\int 3\sqrt{\frac{u-1}{u}}du.$$
This introduces branch cut problems in the integral that we must correct for later, but it allows us to make the substitution $t = \sqrt{(u-1)/u}$, then $dt = 1/2 u^{-2} ((u-1)/u)^{-1/2} du$ and we have 
$$\int \frac{6t^2}{\left(t^2-1\right)^2}dt = -\frac{3t}{t^2-1} - 2 \tan^{-1}(t).$$
When substituting back for $t$, we must use $t = \sqrt{u-1}/\sqrt{u}$ to obtain a correct antiderivative. Finally, we have 
\begin{multline*}
\int x^{-5/6}\sqrt{\sqrt[3]{x} - 1}\, dx =\\ 
3 \sqrt[6]{x}\sqrt{\sqrt[3]{x} - 1} - 3 \tan^{-1}\left(\frac{\sqrt{\sqrt[3]{x} - 1}}{\sqrt[6]{x}}\right).
\end{multline*}


\bigskip

Chebychev's work on the non-existence of algebraic integrals would not be substantially surpassed until the 1970's with the work of Risch\cite{risch1969}\cite{risch1969}, Davenport\cite{davenport1979}, and Trager\cite{trager1984}. \\

\medskip

\problem{2\label{chebychev_1},\pageref{hint:chebychev_1},\pageref{soln:chebychev_1}}
Compute the following Chebychev integrals
\allowdisplaybreaks
\begin{align*}
\text{(i)} & \int \sqrt[3]{x} \left(a\, \sqrt{x}+b\right)^3 dx, \\
\text{(ii)} & \int x^{3/2} \sqrt[3]{a\, x^{5/6}+b} \,dx, \\
\text{(iii)} & \int \frac{x^{4/3}}{\left(a\, \sqrt{x}+b\right)^{2/3}} dx, \\
\text{(iv)} & \int \frac{\left(a\, x^{3/4}+b\right)^{2/3}}{x^3} dx.
\end{align*}

\medskip

\problem{2\label{chebychev_3},\pageref{hint:chebychev_3},\pageref{soln:chebychev_3}}
Use Chebychev's Theorem to determine if the following integrals are elementary. If elementary, compute their antiderivative.
\allowdisplaybreaks
\begin{align*}
\text{(i)} & \int \left(a \sqrt[5]{x}+b\right)^{3/4} \dd x, \\
\text{(ii)} & \int \sqrt[3]{a\, x^2+b} \, \dd x, \\
\text{(iii)} & \int \sqrt[3]{x} \sqrt[3]{a\, x^2+b} \, \dd x, \\
\text{(iv)} & \int \sqrt{a\, x^3+b}\,\dd x, \\
\text{(v)} & \int \frac{\dd x}{\sqrt[3]{a\, x^3+b}}, \\
\text{(vi)} & \int \frac{x\,\dd x}{\sqrt[3]{a\, x^3+b}}, \\
\text{(vii)} & \int \frac{x^4}{\left(a\, x^4+b\right)^{5/4}} \dd x, \\
\text{(viii)} & \int x^{4/5} \left(\frac{a}{x^{4/5}}+b\right)^{5/4} \dd x, \\
\text{(ix)} & \int \frac{x^{3/5}}{\left(a\, x^{3/5}+b\right)^{3/5}} \dd x, \\
\text{(x)} & \int \frac{x^{3/2}}{\sqrt{\frac{a}{\sqrt{x}}+b}} \dd x.
\end{align*}

\medskip

\problem{3\label{chebychev_3b},\pageref{hint:chebychev_3b},\pageref{soln:chebychev_3b}}
For each of the following integrals, transform the integrand into a form suitable for 
applying Chebyshev's Theorem. If an elementary antiderivative exists, compute it.
\allowdisplaybreaks
\begin{align*}
\text{(i)} & \int \sqrt{\sin(x)}\, \dd x, \\
\text{(ii)} & \int \sqrt{\cos(x)}\, \dd x, \\
\text{(iii)} & \int \sqrt{\sinh(x)}\, \dd x, \\
\text{(iv)} & \int \sqrt{\cosh(x)}\, \dd x, \\
\text{(v)} & \int \sqrt{\tan(x)}\, \dd x, \\
\text{(vi)} & \int \sqrt{\tanh(x)}\, \dd x, \\
\text{(vii)} & \int \sqrt{\cot(x)}\, \dd x, \\
\text{(viii)} & \int \sqrt{\coth(x)}\, \dd x.
\end{align*}

\medskip

\problem{3\label{chebychev_3c},\pageref{hint:chebychev_3c},\pageref{soln:chebychev_3c}}
For integer, $n > 1$, transform each integral into a Chebychev integral, if possible, then determine if the integral is elementary. If elementary compute its integral. 
\allowdisplaybreaks
\begin{align*}
\text{(i)} & \int \frac{\dd x}{\sqrt[n]{\sin(x)}}, \\
\text{(ii)} & \int \frac{\dd x}{\sqrt[n]{\cos(x)}}, \\
\text{(iii)} & \int \frac{\dd x}{\sqrt[n]{\tan(x)}}, \\
\text{(iv)} & \int \frac{\dd x}{\sqrt[n]{\csc(x)}}, \\
\text{(v)} & \int \frac{\dd x}{\sqrt[n]{\sec(x)}}, \\
\text{(vi)} & \int \frac{\dd x}{\sqrt[n]{\cot(x)}}.
\end{align*}

\medskip

% ((-1 - 2*x + x^2)*(-1 + (2*Sqrt[-1 + x])/Sqrt[1 + x^2])^(2/3))/((Sqrt[-1 + x]/Sqrt[1 + x^2])^(2/3)*(1 + x^2)^2)
% ((-1 - x - x^2 + x^3)*(-1 + Sqrt[-1 + x]/(Sqrt[1 + x]*Sqrt[1 + x^2]))^(2/3))/(Sqrt[-1 + x]*(1 + x)^(3/2)*(Sqrt[-1 + x]/(Sqrt[1 + x]*Sqrt[1 + x^2]))^(2/3)*(1 + x^2)^(3/2))
\problem{4\label{chebychev_2},\pageref{hint:chebychev_2},\pageref{soln:chebychev_2}}
The following integrals may be transformed to Chebychev integrals
\allowdisplaybreaks
\begin{align*}
\text{(i)} & \int \frac{\left(x^2-2 x-1\right) \left(\frac{2\sqrt{x-1}}{\sqrt{x^2+1}}-1\right)^{2/3}}{\left(x^2+1\right)^2 \left(\frac{\sqrt{x-1}}{\sqrt{x^2+1}}\right)^{2/3}} dx, \\
\text{(ii)} & \int \frac{\left(x^3-x^2-x-1\right) \left(\mathcal{Q}(x)-1\right)^{2/3}}{\sqrt{x-1} (x+1)^{3/2} \left(x^2+1\right)^{3/2} \left(\mathcal{Q}(x)\right)^{2/3}} dx,
\end{align*}
where $$\mathcal{Q}(x) = \frac{\sqrt{x-1}}{\sqrt{x+1} \sqrt{x^2+1}}.$$

\medskip

\problem{3\label{chebychev_4},\pageref{hint:chebychev_4},\pageref{soln:chebychev_4}}
Use Chebychev's Theorem to determine if the following integral is elementary
\allowdisplaybreaks
\begin{align*}
\text{(i)} & \int \frac{dx}{\sqrt{x^{2n}\pm1}},\\
\intertext{for integer $n > 1$. Using this result, determine if the following integrals are elementary}
\text{(ii)} & \int \sqrt[n]{\sec(x)}\,dx, \\
\text{(iii)} & \int \sqrt[n]{\csc(x)}\,dx,\\
\text{(iv)} & \int \sqrt[n]{\text{sech}(x)}\,dx,\\
\text{(v)} & \int \sqrt[n]{\text{csch}(x)}\,dx.
\end{align*}

\medskip

\problem{5\label{chebychev_5},\pageref{hint:chebychev_5},\pageref{soln:chebychev_5}}
Based on experiments with FriCAS and \mma\, we have made the following conjecture regarding a generalised form of the Chebychev integral.

\medskip

\noindent\textsf{Conjecture} \qquad \textit{The integral,
$$\int \left(a\,x^m+b\right)^p\left(c\,x^n+d\right)^q\,dx,$$
for all $a,b,c,d \in \mathbb{R}$ and $n,m,p,q \in \mathbb{Q}$ is elementary if and only if} \\

\noindent\textsc{(base) type I} \qquad \textit{$b\,d = 0$, in which case we have a Chebychev integral.}

\medskip

\noindent\textsc{(base) type II} \qquad \textit{$p \in \mathbb{N}$ or $q \in \mathbb{N}$, in which case we expand and integrate a sum of Chebychev integrals.} 

\medskip

\noindent\textsc{type III (square / square root)} 
\begin{itemize}
\item[--] $p,q = 1/2, -1/2$, and $m,n = 2,-2$, or $m,n = -2,2$.
\item[--] $p,q = 1/2, 1/2$ and $m,n = 2,-2$. 
\item[--] $p,q = -1/2, -1/2$ and $m,n = 2,-2$. 
\end{itemize}

\medskip

\noindent\textsc{type IV (reciprocal exponents)} \qquad \textit{$p+q = 0$ and $1/m = 1/n \in \mathbb{Z}$.}

\medskip

\noindent\textsc{type V (integer combinations)} \qquad \textit{$p+q \in \mathbb{Z}$ and $m,n=1,1$, or $m,n=-1,-1$, or $m,n = 1/\text{den}(p), 1/\text{den}(q)$, or $m,n = -1/\text{den}(p), -1/\text{den}(q)$.} \\

We would be very happy to see this conjecture resolved. If the conjecture is proved, then we would have a nice result on the existence and non--existence of elementary integrals, and if its disproved, then we will have more elementary integrals that are unknown to the author. \\

Compute the following \textsc{type III} integrals
\allowdisplaybreaks
\begin{align*}
\text{(i)} & \int \frac{\sqrt{a\, x^2+b}}{\sqrt{\frac{c}{x^2}+d}} dx,\\
\text{(ii)} & \int \frac{\sqrt{\frac{a}{x^2}+b}}{\sqrt{c\, x^2+d}} dx, \\
\text{(iii)} & \int \sqrt{a\, x^2+b} \sqrt{\frac{c}{x^2}+d}\, dx, \\
\text{(iv)} & \int \frac{dx}{\sqrt{a\, x^2+b} \sqrt{\frac{c}{x^2}+d}}.
\end{align*}

\medskip

Compute the following \textsc{type IV} integrals 
\allowdisplaybreaks
\begin{align*}
\text{(v)} & \int \frac{\left(a \sqrt{x}+b\right)^{2/3}}{\left(c \sqrt{x}+d\right)^{2/3}} dx, \\
\text{(vi)} & \int \frac{\sqrt[4]{\frac{a}{\sqrt[3]{x}}+b}}{\sqrt[4]{\frac{c}{\sqrt[3]{x}}+d}} dx.
\end{align*}

\medskip

Compute the following \textsc{type V} integrals
\allowdisplaybreaks
\begin{align*}
\text{(vii)} & \int \left(a\, x^{-1}+b\right)^{1/3} \left(c\, x^{-1}+d\right)^{2/3} dx, \\
\text{(viii)} & \int \frac{\left(a \sqrt[3]{x}+b\right)^{5/3}}{\left(c \sqrt[3]{x}+d\right)^{2/3}} dx.
\end{align*}

\medskip

(ix) Resolve this conjecture. 

\medskip

An important distinction between our conjecture and Chebychev's Theorem concerns the parameters $a,b,c,d$. Our conjecture requires a solution for all $a,b,c,d$. So, for example the integral,
$$\int \frac{\sqrt[n]{a\, x^n+b}}{\sqrt[n]{c\, x^{2 n} + d}}dx,$$
is not elementary for all $a,b,c,d$. However is it elementary for $c=-a^2$ and $d = b^2$. 


















